% Ek B -- VOC ek sonuçları (revizyon-v2: ana metnin VOC bulguları ve mekanizma bölümlerinden taşınan ikincil içerik)
% main.tex ve cevrimici_ek.tex \input eder.
\setcounter{table}{0}\setcounter{figure}{0}
\section{VOC ek sonuçları}\label{sec:ek-voc}

Bu ek, diğer kartların mekanizma notlarını, iki yöntem ayrıntısını (kutu-tabanlı copy-paste adımları; H07 Soft-NMS tanımı ve doğrulama prosedürü) ve ana metindeki VOC bulgularını tamamlayan mutlak değerleri, ikincil ölçütleri, H07 satırlarını, hata türü kırılımını, copy-paste doz karşılaştırmasını ve nitel paneli verir. Buradaki ikincil ölçütlerin hepsi keşifseldir.

\subsection{Diğer kartların mekanizma notları}\label{sec:ek-mekanizma}
Bu alt bölüm ana metinde birer cümleye indirilen mekanizma notlarının tamamıdır (VOC doğrulaması; puan, tohum-eşli $\Delta$, üç tohum ort. $\pm$ s.d.).

\paragraph{Ölçek ekseni: H01 ve H04.}
H01 eğitim girdisini batch başına 480--800 px arasında değiştirdi (üç tohumda da 11 farklı boyut); H04 mosaic ölçek aralığını daralttı (her tohumda 97\,200 batch'in tümü etkilendi). Müdahale iki kartta da devredeydi, ama ikisi de mAP50-95'i yükseltmedi: H01'in tohum işaretleri karışıktır, H04'ün üçü de negatiftir ve aralığı sıfırı kapsar. H04 5 epoch'luk kısa ön denemede (tohum 0) taban modelin +0.56 önündeyken 100 epoch'ta $-$0.85'e döndü (tohum 0: $-$0.30); kısa ön deneme son sıralamayı öngörmeyen bir kod sağlığı kapısıdır (tek kartlık örnek). Genlik büyütülmediği için ``varsayılan ölçek genliği optimumdur'' denemez.

\paragraph{Kutu kaybı: H02 ve H06.}
H02'de (Inner-CIoU, $r=0.75$) kayba giren IoU $-$ Inner-IoU farkı üç tohumda ortalama 0.027'dir; terim devrededir. mAP50-95 üç tohumda da düştü (Holm-7 sonrası $p=0.244$) ve hedeflenen AP$_{75}$ yükselmedi ($-$0.42 $\pm$ 0.27). H06 (kutu kazancı 7.5 $\to$ 10) AP$_{75}$'i de değiştirmedi (+0.07 $\pm$ 0.30). Taban modelin yüzüncü epoch kayıp payları kutu \%30.1, sınıflandırma \%23.5, DFL \%46.3'tür; en büyük pay olan DFL'nin ağırlığı denenmedi. AP$_{75}$ açığı kutu-IoU teriminin ne biçimiyle ne ağırlığıyla kapandı; atama, kalite hedefi ya da DFL aday açıklamalardır. Inner-IoU'nun raporlanan kazançlarının \cite{zhang_2023_inneriou} COCO ön-eğitimli, TAL atamalı bir ince ayara taşınacağı varsayımı desteklenmedi.

\paragraph{Copy-paste dozu: H03 ve H08.}
H03 yapıştırmayı mosaic'lerin \%11.1'ine, H08 \%47.8'ine uyguladı; ölçülen doz sırasıyla 0.142 ve 0.671 nesne/mosaic'tir. H03'ün tohum-0 farkı ön-eleme eşiğini geçti, ama üç tohum ortalaması SÜRDÜR ölçütünün altında kaldı; mAP50 kazancı üç tohumda da pozitifken AP$_{75}$ farkı $-$0.32 $\pm$ 0.18'dir. Doz ve havuz birlikte genişleyince mAP50 kazancı +0.74 $\pm$ 0.16'dan (H03) +0.25 $\pm$ 0.14'e (H08) indi; artan bir etki gösterilemedi (Şekil~\ref{fig:f5-doz}). ``Doygun mekanizma'' yorumu, ayrıştırılmamış iki değişken nedeniyle zayıf bir hipotezdir. Doğrulamadaki mAP50 izi testte görülmez (Bölüm~\ref{sec:bulgu-tutarlilik}); küçük nesne örneklerini çoğaltmak, çözünürlükle ilişkili görünen açığı bu rejimde kapatmadı (Bölüm~\ref{sec:bulgu-gecikme}).

\paragraph{H05 ayrıntıları.}
H05'in doğrulama mAP50-95 kaybı $-$2.87 $\pm$ 0.33 puandır. VFL gradyan normu BCE'nin 0.34'üdür (üç tohumda 0.342, 0.341, 0.341; kayıp değeri 0.031; tanım: makale Ek~\ref{sec:ek-gradyan}). Uygulama iki birim testinden geçti: $\alpha=1$, $\gamma=0$ ve $y=0$ (yalnız negatif hücreler) iken kayıp BCE ile birebir eşleşir (pozitifler hâlâ $t$ ile ağırlıklanır; H05-pos'un tanımı); pozitif dalın $w=t$ ağırlığı ve gradyanı, elle yazılmış bir karşılaştırma uygulamasına karşı çift duyarlıkla denetlendi. İlk koşumu durduran güvenlik eşiği, taban modelin 5. epoch değerinden 0.03 çıkarılarak H02'nin BCE eğrisinden devralınmıştı.

\subsection{Kutu-tabanlı copy-paste adımları}\label{sec:ek-cp}
Dönüşüm (Bölüm~\ref{sec:mudahaleler}) mosaic'ten sonra, RandomPerspective'ten önce çalışır. Kesişimin ilk kutunun alanına oranı (intersection over area, IoA) $\mathrm{IoA}(b_j,b_l)=|b_j\cap b_l|/|b_j|$ olmak üzere, olasılık $p$ ile şu adımlar uygulanır:
\begin{enumerate}
\item Aday kümesi $\mathcal A=\{\,j:\sqrt{A_j}\le\tau_A,\ \max_{l\ne j}\mathrm{IoA}(b_j,b_l)\le0.1\,\}$ kurulur; $A_j$ kutunun mosaic tuvalindeki piksel alanıdır ve IoA koşulu yamaya başka nesne parçası girmesini önler.
\item $k\sim U\{1,2,3\}$ çekilir ve $\mathcal A$'dan $\min(k,|\mathcal A|)$ nesne rastgele seçilir.
\item Her nesne yeniden ölçeklenmeden, RandomPerspective sonrasında kalan merkez $s\times s$ penceresinde, mevcut tüm kutularla (yeni yapıştırılanlar dahil) kesişimi boş bir konuma yerleştirilir; en çok 10 deneme yapılır, yer bulunamazsa aday atlanır.
\item Piksel yaması kenarları 3 piksellik alfa harmanlamasıyla yapıştırılır (fikir: \cite{dwibedi_2017_cutpaste}; 3 piksel bu çalışmanın seçimidir); sınıf etiketi korunur ve kaynak kutunun boyutuyla yeni bir kutu eklenir.
\end{enumerate}

\subsection{H07: Soft-NMS tanımı ve doğrulama prosedürü}\label{sec:ek-h07-dogrulama}
Soft-NMS satırlarında (S) tutulan kutu $m$ ile örtüşen $b_i$'nin skoru sınıf başına Gauss kuralıyla bozulur \cite{bodla_2017_softnms}:
\begin{equation}\label{eq:softnms}
s_i\leftarrow s_i\exp\!\bigl(-\mathrm{IoU}(m,b_i)^2/\sigma\bigr),\qquad\sigma\in\{0.3,\,0.5\};
\end{equation}
aday havuzu ön-NMS'siz, skoru $\ge0.001$ olan en yüksek 1000 tespittir.

Doğrulama kümesi sıralı dosya adına göre çift ve tek indeksle iki yarıya (A ve B, 500/500) bölünür. Aile başına yarım A'da R0'dan $\Delta>0$ veren en iyi satır aday olur ve şu üç koşulu birlikte sağlarsa DOĞRULANDI sayılır: (a) yarım B'de $\Delta\ge+0.30$; (b) tam doğrulama kümesinde görüntü-düzeyi eşli bootstrap'ın ($B=200$) tek yönlü yüzde 90 alt sınırı sıfırın üstünde; (c) taban modelin tohum 1 ve 2 ağırlıklarında aynı satırın ortalama $\Delta$'sı $\ge+0.20$ ve iki değer de pozitif. Gecikme sınıflaması (KAZANÇ-A/B) için bkz. Bölüm~\ref{sec:mudahaleler}.

\subsection{Mutlak değerler ve tohum 1--2 ortalamaları}

% uretici: evidence/scripts/build_tables.py <- evidence/evidence.json (elle duzenleme bir sonraki uretimde kaybolur)
\begin{table}[!htbp]
\centering\small\setlength{\tabcolsep}{3pt}
\caption{VOC mutlak mAP50-95 [puan] (ort. $\pm$ s.d., n=3) ve tohum 1--2 ortalama $\Delta$. Renk: yalnız Holm-anlamlı (kırmızı kötü).}\label{tab:ekb-mutlak}
\begin{tabularx}{\textwidth}{>{\raggedright\arraybackslash}X c c c}
\toprule
Yöntem & Doğrulama mAP50-95 & Test mAP50-95 & $\Delta$ tohum 1--2 \\
\midrule
Taban model & 63.07 $\pm$ 0.41 & 63.49 $\pm$ 0.11 & -- \\
\midrule
H01 multi\_scale & 62.86 $\pm$ 0.20 & -- & $-$0.35 \\
H02 Inner-CIoU & 62.72 $\pm$ 0.32 & -- & $-$0.41 \\
H03 copy-paste & 63.32 $\pm$ 0.17 & 63.37 $\pm$ 0.20 & +0.19 \\
H04 scale & 62.22 $\pm$ 0.41 & -- & $-$1.13 \\
H05 Varifocal & \kotu{60.20 $\pm$ 0.17} & -- & $-$2.86 \\
H06 kutu kazancı & 63.11 $\pm$ 0.55 & -- & +0.27 \\
H08 yüksek doz & 62.99 $\pm$ 0.14 & -- & $-$0.24 \\
\midrule
Taban model + TTA (eğitimsiz) & -- & 64.32 $\pm$ 0.11 & -- \\
\bottomrule
\addlinespace[3pt]
\multicolumn{4}{@{}p{\dimexpr\textwidth-2\tabcolsep\relax}@{}}{$\Delta$ tohum 1--2: ön-eleme tohumu s0 hariç bağımsız ortalama (doğrulama). Test yalnız taban model, H03 ve TTA için açıldı. Renk ölçütü: doğrulamada Holm-7, testte Holm-12 ve sıfırı kesmeyen GA (Tablo~\ref{tab:voc-ana}).} \\
\end{tabularx}
\end{table}


Tablo~\ref{tab:ekb-mutlak} doğrulama ve test mAP50-95'in mutlak değerlerini ve ön-eleme tohumu hariç tohum 1--2 ortalamasını verir. Taban model doğrulamada 63.07 $\pm$ 0.41, testte 63.49 $\pm$ 0.11 puandır; TTA'lı test değeri 64.32 $\pm$ 0.11'dir. Elenen kartların bağımsız tohum 1--2 ortalamalarından hiçbiri $+0.30$'u aşmaz; en yakını H06'dır ($+0.27$; tohum 0'da $-$0.42). Tohum 0'da elenen H06'nın bağımsız tohumlarda eşiğe en yakın kart olması, tek tohumlu kapının yanlış-negatif riski taşıdığını hatırlatır (Bölüm~\ref{sec:bulgu-ablasyon}).

\subsection{Test ikincil ölçütleri}

% uretici: evidence/scripts/build_tables.py <- evidence/evidence.json (elle duzenleme bir sonraki uretimde kaybolur)
\begin{table}[!htbp]
\centering\small\setlength{\tabcolsep}{3pt}
\caption{VOC ikincil ölçütler: taban modele göre tohum-eşli $\Delta$ [puan], n=3. $p$: çift yönlü eşli $t$; Holm-12: 12 testlik aile.}\label{tab:ekb-test-ikincil}
\begin{tabularx}{\textwidth}{>{\raggedright\arraybackslash}X l c c c}
\toprule
Karşılaştırma & Metrik & $\Delta$ ort.$\pm$s.d. & $p$ & Holm-12 $p$ \\
\midrule
H03 kutu copy-paste (test) & mAP50 & $-$0.02 $\pm$ 0.21 & 0.883 & 1.000 \\
 & mAP75 & $-$0.11 $\pm$ 0.13 & 0.268 & 1.000 \\
 & AP$_S$ & +0.81 $\pm$ 2.55 & 0.639 & 1.000 \\
 & AP$_M$ & +0.48 $\pm$ 0.14 & 0.027$^\dagger$ & 0.222 \\
 & AP$_L$ & $-$0.23 $\pm$ 0.09 & 0.048$^\dagger$ & 0.333 \\
\midrule
Taban model + TTA (test) & mAP50 & +0.10 $\pm$ 0.27 & 0.579 & 1.000 \\
 & mAP75 & +1.12 $\pm$ 0.30 & 0.023$^\dagger$ & 0.222 \\
 & AP$_S$ & $-$0.48 $\pm$ 0.82 & 0.416 & 1.000 \\
 & AP$_M$ & $-$1.07 $\pm$ 0.18 & 0.009$^\dagger$ & 0.112 \\
 & AP$_L$ & +1.50 $\pm$ 0.36 & 0.019$^\dagger$ & 0.210 \\
\midrule
H03 (doğrulama) & mAP50 & +0.74 $\pm$ 0.16 & 0.015 & -- \\
\bottomrule
\addlinespace[3pt]
\multicolumn{5}{@{}p{\dimexpr\textwidth-2\tabcolsep\relax}@{}}{mAP50-95 satırları: Tablo~\ref{tab:voc-ana}. Holm-12: K5 ailesi (H03, TTA $\times$ 6 metrik; H03 mAP50-95 dahil, konservatif); en küçüğü 0.112, renk yok. $^\dagger$: düzeltmesiz $p<0.05$, Holm sonrası anlamlı değil.} \\
\end{tabularx}
\end{table}


Tablo~\ref{tab:ekb-test-ikincil} H03 ve TTA'nın test ikincil ölçütlerini verir. Düzeltmesiz $p<0.05$ veren farklar vardır (H03'te AP$_M$ ve AP$_L$; TTA'da mAP75, AP$_M$ ve AP$_L$), ancak on iki testlik ailede Holm sonrası en küçük $p$ 0.112'dir; bu farklar yön olarak okunur. TTA'nın test mAP50-95 farkı için tohum düzeyindeki $t$ aralığı [+0.29, +1.38]'dir.

H03'ün doğrulamada üç tohumda pozitif olan mAP50 kazancı (+0.74 $\pm$ 0.16, $p=0.015$) testte $-$0.02 $\pm$ 0.21'dir. Doğrulamanın \%70'i VOC2012 olduğundan kaynak farkı sınandı: mAP50 farkı VOC07 alt kümesinde +0.73 $\pm$ 1.23, VOC12'de +0.64 $\pm$ 0.58'dir (303/697 görüntü). Hipotez desteklenmedi ve \texttt{best.pt} seçim iyimserliği sınanmadı.

\subsection{H07 eğitimsiz çıkarım satırları}

% uretici: evidence/scripts/build_tables.py <- evidence/evidence.json (elle duzenleme bir sonraki uretimde kaybolur)
\begin{table}[!htbp]
\centering\small\setlength{\tabcolsep}{3pt}
\caption{H07 eğitimsiz çıkarım satırları (doğrulama, tohum 0) ve TTA test satırı: $\Delta$ mAP50-95 ve gecikme. Mavi doğrulandı; kırmızı $\Delta\le-0.5$ ya da $>$1.10$\times$.}\label{tab:h07}
\begin{tabularx}{\textwidth}{l >{\raggedright\arraybackslash}X c c c c c l}
\toprule
Satır & Tanım & $\Delta$ mAP50-95 & $\Delta$ AP$_S$ & $\Delta$ AP$_L$ & [ms] & [$\times$] & Durum \\
\midrule
R0 & Referans (NMS 0.7, 640 px) & 0.00 & 0.00 & 0.00 & 8.45 & 1.00 & referans \\
N50 & NMS IoU 0.5 & \kotu{$-$0.54} & $-$0.93 & $-$0.56 & 8.25 & 0.98 & aday değil \\
N60 & NMS IoU 0.6 & $-$0.21 & $-$0.20 & $-$0.20 & 8.30 & 0.98 & aday değil \\
N80 & NMS IoU 0.8 & $-$0.10 & +0.03 & $-$0.03 & 8.33 & 0.99 & aday değil \\
S03 & Soft-NMS $\sigma$=0.3 & $-$0.12 & $-$0.11 & $-$0.03 & 9.10 & 1.08 & aday değil \\
S05 & Soft-NMS $\sigma$=0.5 & \kotu{$-$0.58} & +0.06 & $-$0.51 & 9.14 & 1.08 & aday değil \\
U704 & test imgsz 704 & $-$0.31 & +2.41 & $-$0.82 & 8.92 & 1.06 & aday değil \\
U768 & test imgsz 768 & \kotu{$-$2.06} & +4.55 & $-$3.41 & 9.43 & 1.12$^\dagger$ & aday değil \\
U832 & test imgsz 832 & \kotu{$-$5.19} & +4.94 & $-$7.24 & 9.61 & 1.14$^\dagger$ & aday değil \\
T1 & TTA 1, 0.83, 0.67 + flip & \iyi{+0.80} & $-$0.12 & +1.58 & 21.68 & \kotu{2.57} & KAZANÇ-B \\
T2 & TTA 1, 1.2, 1.4 + flip & \kotu{$-$4.12} & +2.75 & $-$5.99 & 21.54 & \kotu{2.55} & aday değil \\
W1 & WBF (T1 dalları) & $-$0.28 & +0.41 & +0.65 & 25.27 & \kotu{2.99} & aday değil \\
W2 & WBF (T2 dalları) & \kotu{$-$5.74} & +5.10 & $-$7.85 & 26.14 & \kotu{3.10} & aday değil \\
\midrule
T1 test & TTA, 3 tohum, test & +0.83 $\pm$ 0.22$^\dagger$ & $-$0.48 & +1.50 & 21.64 & \kotu{2.29} & KAZANÇ-B \\
\bottomrule
\addlinespace[3pt]
\multicolumn{8}{@{}p{\dimexpr\textwidth-2\tabcolsep\relax}@{}}{Doğrulama: taban model tohum-0 \texttt{best.pt} (R0 = 8.45 ms); TTA test: kendi oturumu (taban model 9.43 ms). Kırmızı $\Delta$ iki doğrulama yarısında da negatif (GA yok). $^\dagger$: oran tek-koşum gürültüsü (1.14$\times$) içinde (renksiz) ya da $\Delta$ düzeltmesiz $p<0.05$.} \\
\end{tabularx}
\end{table}


Tablo~\ref{tab:h07}, H07'nin on iki eğitimsiz satırını R0 referansına göre ve TTA test satırını gecikmeyle verir; Şekil~\ref{fig:f3-pareto} aynı satırları çizer (Bölüm~\ref{sec:bulgu-gecikme}). T1 tek DOĞRULANDI satırıdır ve gecikme oranı nedeniyle KAZANÇ-B sayılır. WBF tanısının iki ek ayrıntısı şudur. WBF'nin IoU eşiği NMS'inkinden gevşek olduğundan 0.6 eşiğinde tek dal bile ham NMS'ten ayrılır. Varsayılan WBF ayarı yüksek IoU'da da kaybettirdi (T1'e göre AP$_{75}$ $-$1.33, AP$_{90}$ $-$2.15; Bölüm~\ref{sec:basarisiz-wbf}).

\subsection{Eğitimsiz çıkarım: WBF tanısı}\label{sec:basarisiz-wbf}
Varsayılan WBF ayarı T1'in mAP50-95'ini $-$1.08 düşürdüğünden doğrulamada, test sonrası ve keşifsel bir tanı yapıldı (kazanç iddiası değildir). Uygulama doğruydu: 8 birim testinin 8'i geçti ve tek dallı WBF, IoU eşiği $\approx$1 iken ham NMS'ten $+0.000$ puan ayrıldı. Kayıp iki ayardan gelir. Ortalama skor füzyonunda üç görünümden azında görülen kutunun skoru küçülür; ölçek dalları her boyuttaki nesneyi göremediğinden bu ceza sıralamayı bozar ve maksimum füzyonu +0.98 puan getirdi. IoU eşiğini 0.7'ye çıkarmak ortalama füzyonda da kaybı azalttı; iki ayarın payları ayrıştırılmadı. Izgaranın en iyi ayarı (IoU 0.7, maksimum) T1'i +0.45 geçti; doğrulamada seçildiği için bu kanıt değil, etkinin ayara bağlı olduğunun göstergesidir.

\subsection{Kart başına ikincil ölçütler (doğrulama)}
Eğitim kartlarının boyut ve eşik ölçütlerindeki farklar şunlardır. H01'de AP$_S$ $-$0.30 $\pm$ 2.24, AP$_M$ +0.84; H04'te AP$_S$ $-$1.06 $\pm$ 0.68, AP$_M$ $-$0.75'tir. H02'de AP$_S$ değişimi +1.39 $\pm$ 1.09'dur. H06'da mAP50 +0.30 $\pm$ 0.55'tir. Bu AP$_S$ değişimlerinin hepsi taban modelin tohum aralığı (2.87 puan) içinde kalır. H05'in kaybı üç boyut sınıfında görülür: AP$_S$ $-$2.39, AP$_M$ $-$4.06, AP$_L$ $-$2.10.

H05'in son 10 epoch'taki doğrulama kazancı +0.20, +0.27, +0.23, taban modelinki +0.04, +0.16, $-$0.03 puandır. Aradaki fark ($\approx$0.18 puan) tek bir doğrulama ölçümünün gürültüsünün (Bölüm~\ref{sec:bulgu-tutarlilik}) çok altındadır ve bu evre mosaic'in kapandığı evredir; yakınsama hipotezi bu yüzden yalnız zayıf desteklenir.

H05c'nin tohum başına değerleri şunlardır: $\Delta_{\mathrm{c},s}$ +0.02, $-$1.15, $-$0.40; $D_s$ +2.90, +2.04, +2.13. Tohum başına geri kazanım 1.01, 0.64 ve 0.84'tür (ortalama 0.83). Kalan farkta iki tohum negatif, biri $\approx$0'dır. H05c koşumlarında ölçülen VFL/BCE gradyan normu oranı 0.347--0.349'dur. En iyi epoch H05c'de 96, 96 ve 95, H05'te 100, 100 ve 100'dür. Karar kuralı, geçersiz bir koşumun yeniden koşulmayıp ``koşulamadı'' olarak raporlanmasını öngörüyordu; büyüklüğün kaybı açıklamadığı durum ise biçim hipotezini güçlendirecek ama kanıtlamayacaktı.

\subsection{Hata türü kırılımı}\label{sec:bulgu-hata}
Tespitler TIDE'a benzer biçimde hata türlerine ayrılır \cite{bolya_2020_tide}; eşleştirme pycocotools ile (en çok 300 tespit) yapılır. Doğru pozitif (true positive, TP), IoU$\ge0.5$ ile eşleşen tespittir; Dupe, GT'si daha yüksek skorlu bir kutuyla alınmış IoU$\ge0.5$ tespit; Cls, başka sınıf GT ile IoU$\ge0.5$; Loc, aynı sınıf GT ile $0.1\le$IoU$<0.5$; Both, başka sınıf GT ile $0.1\le$IoU$<0.5$; Bkg, hiçbir GT ile IoU$\ge0.1$ olmayan tespit; Miss, hiçbir kutuyla IoU$\ge0.5$ eşleşmeyen GT'dir. Oracle $\Delta$AP$_{50}$, ilgili hata düzeltilip (Loc: kutu GT'ye yapıştırılır, Cls: sınıf düzeltilir, Dupe, Both ve Bkg: kutu silinir, Miss: GT silinir) yeniden değerlendirildiğinde AP$_{50}$'deki artıştır. Kırılım işletim noktasında (skor $\ge0.25$) ve eşiksiz ($\ge0.001$), boyut ve sınıf düzeyinde ayrıca raporlanır.

% F2 -- hata turu kirilimi (TIDE-benzeri oracle ΔAP50): (a) val, (b) test. Veri: figures/data/f2a_val.dat, f2b_test.dat
% (uretici: figures/scripts/f2_hata_turu.py <- evidence/evidence.json). Figur ici metin Turkce.
\begin{figure}[!htbp]
\centering
\pgfplotstableset{col sep=semicolon}
\pgfplotsset{
  ftwo/.style={scale only axis, width=0.86\textwidth, height=4.3cm, ybar, bar width=8pt, symbolic x coords={Cls,Loc,Both,Dupe,Bkg,Miss}, xtick=data,
    enlarge x limits=0.10, ymin=0, ymax=4.6, ytick={0,1,2,3,4}, font=\small, tick label style={font=\small}, label style={font=\small},
    title style={font=\small, yshift=-1pt}, ymajorgrids=true, grid style={black!10}, axis line style={black!70},
    ylabel={Oracle $\Delta$AP50 [puan]}, error bars/y dir=both, error bars/y explicit,
    error bars/error bar style={line width=0.9pt, black}, error bars/error mark options={mark size=2pt, black}},
  fbase/.style={fill=black!65, draw=black!65}, fhthree/.style={fill=iyiRenk, draw=iyiRenk}, fhsix/.style={fill=kotuRenk, draw=kotuRenk},
}
\begin{tikzpicture}
% ---------------- (a) validation
\begin{axis}[name=A, ftwo, title={(a) Doğrulama (1000 görüntü; taban model ve H03: 3 tohum, H06: 1 tohum)},
  legend style={at={(0.5,1.28)}, anchor=south, legend columns=3, font=\footnotesize, draw=none}, legend image code/.code={\draw[#1] (0cm,-0.1cm) rectangle (0.3cm,0.1cm);}]
\addplot[fbase] table[x=tur, y=baseline_mean, y error plus=baseline_ep, y error minus=baseline_em]{\figroot data/f2a_val.dat};
\addplot[fhthree] table[x=tur, y=h03_mean, y error plus=h03_ep, y error minus=h03_em]{\figroot data/f2a_val.dat};
\addplot[fhsix] table[x=tur, y=h06_mean, y error plus=h06_ep, y error minus=h06_em]{\figroot data/f2a_val.dat};
\legend{Taban model, H03 kutu copy-paste, H06 kutu kazancı 10}
\end{axis}
% ---------------- (b) test
\begin{axis}[name=B, ftwo, at={(A.south west)}, anchor=north west, yshift=-1.6cm, title={(b) Test (4952 görüntü; taban model ve H03: 3 tohum)},
  xlabel={Oracle ile düzeltilen hata türü}]
\addplot[fbase] table[x=tur, y=baseline_mean, y error plus=baseline_ep, y error minus=baseline_em]{\figroot data/f2b_test.dat};
\addplot[fhthree] table[x=tur, y=h03_mean, y error plus=h03_ep, y error minus=h03_em]{\figroot data/f2b_test.dat};
\end{axis}
\end{tikzpicture}
\caption{Hata türü oracle'ı: ilgili hata düzeltilirse $\Delta$AP50 [puan] (Cls sınıf, Loc konum, Both ikisi, Dupe yinelenen, Bkg arka plan, Miss kaçırılan). Çubuk ortalama, bıyık tohum min–maks; (a) doğrulama, (b) test.}\label{fig:f2-hata-turu}
\end{figure}


Şekil~\ref{fig:f2-hata-turu} hata türlerinin oracle $\Delta$AP50 maliyetini taban model, H03 ve H06 için karşılaştırır (test kırılımı keşifsel; H06 yalnız valda, tohum 0). Arka plan yanlış pozitifi (Bkg) iki kümede de en maliyetli türdür (val 3.48, test 3.79 puan; Miss: 3.32 ve 2.03). H03 ve H06 çubukları taban modelden valda en çok 0.31, testte en çok 0.06 puan ayrılır; hata profili kaymadı ve bu, mAP'teki sıfıra yakın farkla tutarlıdır.

\subsection{Copy-paste doz karşılaştırması}

% F5 -- doz-yanit: copy-paste dozu (nesne/mosaic) vs Δ mAP50-95 ve Δ mAP50 (val, n=3). Veri: figures/data/f5_doz.dat, f5_ort.dat, f5_ref.dat
% (uretici: figures/scripts/f5_doz.py <- evidence/evidence.json). Figur ici metin Turkce. xtick konumlari f5_ref.dat'tan okunur (gomulu sabit yok).
% Not: H08 olasiligi (0.5->1.0) ve aday havuzunu (<=48 -> <=96 px) birlikte degistirir; doz ile havuz ayrilmamistir.
\begin{figure}[!htbp]
\centering
\pgfplotstableset{col sep=semicolon}
\pgfplotstableread{\figroot data/f5_ref.dat}\ffiveref
\pgfplotstablegetelem{0}{sd5095}\of\ffiveref \edef\fSdA{\pgfplotsretval}
\pgfplotstablegetelem{0}{sd50}\of\ffiveref \edef\fSdB{\pgfplotsretval}
\pgfplotstablegetelem{0}{doz_lo}\of\ffiveref \edef\fDozLo{\pgfplotsretval}
\pgfplotstablegetelem{0}{doz_h03}\of\ffiveref \edef\fDozA{\pgfplotsretval}
\pgfplotstablegetelem{0}{doz_h08}\of\ffiveref \edef\fDozB{\pgfplotsretval}
\pgfplotsset{
  ffive/.style={scale only axis, width=0.80\textwidth, height=3.6cm, xmin=-0.05, xmax=0.95, xtick={0,\fDozA,\fDozB,\fDozLo},
    xticklabel style={/pgf/number format/fixed, /pgf/number format/precision=3}, font=\small, tick label style={font=\small}, label style={font=\small},
    title style={font=\small, yshift=-1pt}, ymajorgrids=true, grid style={black!10}, axis line style={black!70}, clip=false},
}
\begin{tikzpicture}
% ---------------- (a) Δ mAP50-95
\begin{axis}[name=A, ffive, ymin=-0.9, ymax=1.2, ytick={-0.5,0,0.5,1}, xticklabels={}, title={(a) $\Delta$ mAP50-95, copy-paste dozuna karşı (doğrulama, 3 tohum)},
  ylabel={$\Delta$ mAP50-95 [puan]}]
\fill[black!13] (axis cs:-0.05,-\fSdA) rectangle (axis cs:0.95,\fSdA);
\fill[iyiRenk!8] (axis cs:\fDozLo,-0.9) rectangle (axis cs:0.95,1.2);
\draw[line width=1.1pt, black] (axis cs:-0.05,0) -- (axis cs:0.95,0);
\node[font=\scriptsize, anchor=north east, align=right, inner sep=1.5pt] at (axis cs:0.95,1.2) {önceden yazılmış H08 hedefi $D\ge0.80$\\(beklenen aralık 0.8--1.2; ulaşılamadı)};
\addplot[mark=none, dashed, line width=1.2pt, black!65, forget plot] table[x=doz, y=d5095]{\figroot data/f5_ort.dat};
\addplot[only marks, mark=*, mark size=2.4pt, black!65, forget plot] table[x=doz, y=d5095]{\figroot data/f5_doz.dat};
\addplot[only marks, mark=|, mark size=6pt, line width=1.6pt, black, forget plot, point meta=explicit symbolic, nodes near coords={\pgfplotspointmeta},
  every node near coord/.append style={font=\scriptsize, anchor=south west, xshift=3pt, yshift=1pt, inner sep=1pt, fill=white}] table[x=doz, y=d5095, meta=lab5095]{\figroot data/f5_ort.dat};
\end{axis}
% ---------------- (b) Δ mAP50
\begin{axis}[name=B, ffive, at={(A.south west)}, anchor=north west, yshift=-1.4cm, ymin=-0.3, ymax=1.2, ytick={0,0.5,1},
  title={(b) $\Delta$ mAP50, copy-paste dozuna karşı (doğrulama, 3 tohum)}, ylabel={$\Delta$ mAP50 [puan]}, xlabel={Doz [mosaic başına yapıştırılan nesne]},
  legend style={legend cell align=left, at={(0.5,-0.55)}, anchor=north, legend columns=2, font=\footnotesize, draw=none}]
\fill[black!13] (axis cs:-0.05,-\fSdB) rectangle (axis cs:0.95,\fSdB);
\fill[iyiRenk!8] (axis cs:\fDozLo,-0.3) rectangle (axis cs:0.95,1.2);
\draw[line width=1.1pt, black] (axis cs:-0.05,0) -- (axis cs:0.95,0);
\addplot[mark=none, dashed, line width=1.2pt, black!65, forget plot] table[x=doz, y=d50]{\figroot data/f5_ort.dat};
\addplot[only marks, mark=*, mark size=2.4pt, black!65, forget plot] table[x=doz, y=d50]{\figroot data/f5_doz.dat};
\addplot[only marks, mark=|, mark size=6pt, line width=1.6pt, black, forget plot, point meta=explicit symbolic, nodes near coords={\pgfplotspointmeta},
  every node near coord/.append style={font=\scriptsize, anchor=south west, xshift=3pt, yshift=1pt, inner sep=1pt, fill=white}] table[x=doz, y=d50, meta=lab50]{\figroot data/f5_ort.dat};
\addlegendimage{only marks, mark=*, mark size=2.4pt, black!65}\addlegendentry{tohum değeri (her dozda 3 tohum)}
\addlegendimage{only marks, mark=|, mark size=6pt, line width=1.6pt, black}\addlegendentry{ortalama}
\addlegendimage{area legend, fill=black!13, draw=none}\addlegendentry{taban model doğrulama $\pm1$ s.d. (3 tohum)}
\addlegendimage{dashed, line width=1.2pt, black!65}\addlegendentry{ortalamaları birleştirir (uydurulmuş eğri değil)}
\end{axis}
\end{tikzpicture}
\caption{Kutu-tabanlı copy-paste dozuna (mosaic başına yapıştırılan nesne) karşı doğrulama $\Delta$ mAP50-95 (a) ve $\Delta$ mAP50 (b), 3 tohum. H08 yapıştırma olasılığını ve aday havuzunu birlikte değiştirir; doz tek başına ayrıştırılmamıştır.}\label{fig:f5-doz}
\end{figure}


Şekil~\ref{fig:f5-doz} doğrulama kazancını yapıştırma dozuna (mosaic başına nesne) karşı çizer: 0 (taban model), 0.142 (H03), 0.671 (H08; H03'ün 4.7 katı). H08 olasılığı ve aday havuzunu birlikte değiştirdiğinden (Bölüm~\ref{sec:mudahaleler}) doz ile boyut karışımı ayrışmaz; bu bir doz--yanıt eğrisi değil, üç noktalı bir karşılaştırmadır. $\Delta$mAP50 +0.74 $\pm$ 0.16'dan +0.25 $\pm$ 0.14'e iner; H08'de $\Delta$AP$_S$ +0.31 $\pm$ 2.84'tür. Doz ve havuz birlikte genişleyince kazanç büyümedi (``artan bir etki gösterilemedi'').

\subsection{Nitel panel}\label{sec:bulgu-nitel}

% F6 -- NITEL PANEL (figure-style A2): VOC07 test, 6 goruntu x {GT, baseline s0, H03 s0, baseline+TTA s0}. Karolar: figures/tiles/f6/r<k>_{gt,base,h03,tta}.jpg
% (gercek cikti karolari; uretici figures/scripts/f6_nitel.py, secim olcutu orada ve caption'da). Etiket/not/lejant TikZ'te (Turkce).
\begin{figure}[!htbp]
\centering
\pgfplotstableset{col sep=semicolon}
\pgfplotstableread{\figroot data/f6_panel.dat}\fsixtbl
\pgfmathsetmacro{\fsixlab}{3.0}% satir etiketi sutunu [cm] (yerlesim)
\pgfmathsetmacro{\fsixw}{(\textwidth/1cm-\fsixlab-0.2)/4-0.10}% karo genisligi [cm], \textwidth'e kilitli
\pgfmathsetmacro{\fsixh}{0.75*\fsixw}% karo yuksekligi [cm] (4:3)
\pgfmathsetmacro{\fsixrow}{\fsixh+0.93}% satir yuksekligi (karo + not)
\begin{tikzpicture}[x=1cm,y=1cm]
% sutun basliklari
\foreach \cn/\t in {0/{GT (gerçek)},1/{Taban model (s0)},2/{H03 (s0)},3/{Taban model + TTA (s0)}}{
  \node[font=\footnotesize\bfseries, anchor=south, text width=\fsixw cm, align=center, inner sep=1pt] at (\fsixlab+\cn*\fsixw+\cn*0.10+0.5*\fsixw, 0.04) {\t};}
\foreach \r in {0,...,5}{
  \pgfmathtruncatemacro{\rn}{\r+1}
  \pgfplotstablegetelem{\r}{kriter}\of\fsixtbl \edef\kr{\pgfplotsretval}
  \pgfplotstablegetelem{\r}{goruntu}\of\fsixtbl \edef\gimg{\pgfplotsretval}
  \pgfplotstablegetelem{\r}{n_gt}\of\fsixtbl \edef\ngt{\pgfplotsretval}
  \pgfplotstablegetelem{\r}{n_kucuk}\of\fsixtbl \edef\nks{\pgfplotsretval}
  \pgfplotstablegetelem{\r}{gt_siniflar}\of\fsixtbl \edef\gcls{\pgfplotsretval}
  \pgfplotstablegetelem{\r}{ap_base}\of\fsixtbl \edef\apb{\pgfplotsretval}
  \pgfplotstablegetelem{\r}{ap_h03}\of\fsixtbl \edef\aph{\pgfplotsretval}
  \pgfplotstablegetelem{\r}{ap_tta}\of\fsixtbl \edef\apt{\pgfplotsretval}
  \pgfplotstablegetelem{\r}{tp_base}\of\fsixtbl \edef\tpb{\pgfplotsretval}
  \pgfplotstablegetelem{\r}{fp_base}\of\fsixtbl \edef\fpb{\pgfplotsretval}
  \pgfplotstablegetelem{\r}{fn_base}\of\fsixtbl \edef\fnb{\pgfplotsretval}
  \pgfplotstablegetelem{\r}{tp_h03}\of\fsixtbl \edef\tph{\pgfplotsretval}
  \pgfplotstablegetelem{\r}{fp_h03}\of\fsixtbl \edef\fph{\pgfplotsretval}
  \pgfplotstablegetelem{\r}{fn_h03}\of\fsixtbl \edef\fnh{\pgfplotsretval}
  \pgfplotstablegetelem{\r}{tp_tta}\of\fsixtbl \edef\tpt{\pgfplotsretval}
  \pgfplotstablegetelem{\r}{fp_tta}\of\fsixtbl \edef\fpt{\pgfplotsretval}
  \pgfplotstablegetelem{\r}{fn_tta}\of\fsixtbl \edef\fnt{\pgfplotsretval}
  \pgfmathsetmacro{\yt}{-\r*\fsixrow}
  % renk: baseline'a gore goruntu-AP'si (iyi mavi kalin / kotu kirmizi egik; esit ise duz)
  \pgfmathparse{\aph>\apb+0.0005}\ifnum\pgfmathresult>0 \def\sth{\iyi}\else\pgfmathparse{\aph<\apb-0.0005}\ifnum\pgfmathresult>0 \def\sth{\kotu}\else\def\sth{\relax}\fi\fi
  \pgfmathparse{\apt>\apb+0.0005}\ifnum\pgfmathresult>0 \def\stt{\iyi}\else\pgfmathparse{\apt<\apb-0.0005}\ifnum\pgfmathresult>0 \def\stt{\kotu}\else\def\stt{\relax}\fi\fi
  % satir etiketi: olcut, goruntu, GT siniflari
  \node[font=\scriptsize, text width=\fsixlab cm-0.2cm, align=left, anchor=north west, inner sep=1pt] at (0,\yt) {\textbf{\kr}\\ görüntü \gimg\\ GT: \gcls\\ (\ngt\ nesne, \nks\ küçük)};
  % karolar
  \node[inner sep=0, anchor=north west] at (\fsixlab,\yt) {\includegraphics[width=\fsixw cm]{\figroot tiles/f6/r\rn_gt.jpg}};
  \node[inner sep=0, anchor=north west] at (\fsixlab+\fsixw+0.10,\yt) {\includegraphics[width=\fsixw cm]{\figroot tiles/f6/r\rn_base.jpg}};
  \node[inner sep=0, anchor=north west] at (\fsixlab+2*\fsixw+0.20,\yt) {\includegraphics[width=\fsixw cm]{\figroot tiles/f6/r\rn_h03.jpg}};
  \node[inner sep=0, anchor=north west] at (\fsixlab+3*\fsixw+0.30,\yt) {\includegraphics[width=\fsixw cm]{\figroot tiles/f6/r\rn_tta.jpg}};
  % notlar (AP = goruntu-duzeyi AP50-95; TP/FP/FN skor>=0.25)
  \node[font=\scriptsize, text width=\fsixw cm, align=center, inner sep=1pt, anchor=north] at (\fsixlab+1.5*\fsixw+0.10,\yt-\fsixh) {AP \apb\\ \tpb/\fpb/\fnb};
  \node[font=\scriptsize, text width=\fsixw cm, align=center, inner sep=1pt, anchor=north] at (\fsixlab+2.5*\fsixw+0.20,\yt-\fsixh) {\sth{AP \aph}\\ \tph/\fph/\fnh};
  \node[font=\scriptsize, text width=\fsixw cm, align=center, inner sep=1pt, anchor=north] at (\fsixlab+3.5*\fsixw+0.30,\yt-\fsixh) {\stt{AP \apt}\\ \tpt/\fpt/\fnt};
}
% lejant
\pgfmathsetmacro{\yleg}{-6*\fsixrow-0.02}
\node[font=\scriptsize, anchor=north west, inner sep=1pt] at (0,\yleg) {Skor $\ge 0.25$ tespitler. Karo altı: görüntü AP50-95; TP/FP/FN sayıları.};
\draw[line width=1.8pt, iyiRenk] (0.1,\yleg-0.62) -- (0.6,\yleg-0.62); \node[font=\scriptsize, anchor=west] at (0.65,\yleg-0.62) {TP (doğru)};
\draw[line width=1.8pt, kotuRenk] (2.65,\yleg-0.62) -- (3.15,\yleg-0.62); \node[font=\scriptsize, anchor=west] at (3.2,\yleg-0.62) {FP (yanlış pozitif)};
\draw[line width=1.8pt, orange] (6.6,\yleg-0.62) -- (7.1,\yleg-0.62); \node[font=\scriptsize, anchor=west] at (7.15,\yleg-0.62) {FN (kaçırılan)};
\draw[line width=2.4pt, white] (9.7,\yleg-0.62) -- (10.2,\yleg-0.62); \draw[line width=0.5pt, black] (9.7,\yleg-0.58) -- (10.2,\yleg-0.58); \draw[line width=0.5pt, black] (9.7,\yleg-0.66) -- (10.2,\yleg-0.66); \node[font=\scriptsize, anchor=west] at (10.25,\yleg-0.62) {gerçek etiket};
\end{tikzpicture}
\caption{VOC07 testinde altı görüntüde GT, taban model, H03 ve TTA tespitleri (tohum 0). Mavi kalın/kırmızı eğik: taban modele göre iyi/kötü. Seçim: $\ge$2 GT'li, üç yöntemde tespitli, taban model AP'si Q1--Q3 içindeki adaylardan TTA kazancı/kaybı, küçük nesneli H03 kazancı/kaybı, rastgele 1.}\label{fig:f6-nitel}
\end{figure}


Şekil~\ref{fig:f6-nitel} aynı test görüntülerinde GT, taban model, H03 ve taban model+TTA tespitlerini (tohum 0, skor $\ge0.25$) iki yönün uç örnekleri ve bir rastgele kareyle gösterir; seçim ölçütü başlıktadır (1642 görüntülük havuz; taban model görüntü-AP Q1--Q3: 0.497--0.896). TTA bir sahnede görüntü-AP'sini 0.522'den 0.971'e çıkarırken bir otomobil sahnesinde 0.895'ten 0.248'e düşürür; H03 de ilk sahnede sıçrar, yani farklar TTA'ya özgü değildir. Tüm geçerli görüntülerde (tohum 0) TTA'nın ortalama görüntü-AP farkı $-$1.26'dır; görüntülerin \%31.8'i iyileşir, \%38.4'ü kötüleşir (H03: +0.36; \%35.5 ve \%34.7). Veri kümesi mAP'i ise +0.88 artar; kazanç tek tip bir görsel iyileşme değildir (keşifsel).

\FloatBarrier
